\section{Emerging Multimodal Architectures}

\begin{frame}{}
    \LARGE Next-Generation Architectures: \textbf{Multimodal Foundation Models}

    \vspace{1em}
    \normalsize What changes when the model must also see, hear and draw?
\end{frame}

\begin{frame}{2026: The Separate Vision Model Disappears}
    \begin{center}
    \small
    \renewcommand{\arraystretch}{1.25}
    \begin{adjustbox}{max width=\linewidth}\begin{tabular}{>{\raggedright\arraybackslash}p{2.8cm}>{\raggedright\arraybackslash}p{3.4cm}>{\raggedright\arraybackslash}p{6.9cm}}
        \toprule
        \textbf{Model} & \textbf{Backbone} & \textbf{How vision is built in} \\
        \midrule
        Kimi K2.5 (Jan 2026) & MoE, 1.04T / 32B active & a 400M native-resolution ViT, with 15T tokens of joint vision and text pre-training \\
        Qwen3.5 (Feb 2026) & 3:1 hybrid MoE, 397B / 17B & one checkpoint replaces the separate Qwen3-VL line \\
        Gemma 4 (Jul 2026) & 2.3B to 31B, dense and MoE & a 550M ViT in most sizes \\
        \bottomrule
    \end{tabular}\end{adjustbox}
    \end{center}
    \begin{itemize}\itemsep0.5em
        \item Until 2025, most labs shipped a text model and a separate ``-VL'' variant. Each of these three ships \textbf{one natively multimodal checkpoint}.
        \item Note the backbones. They are the hybrid and MoE designs from earlier today. The novelty is in the backbone and in token budgets, not in a new fusion trick.
    \end{itemize}
    \src{Kimi Team, ``Kimi K2.5,'' \href{https://arxiv.org/abs/2602.02276}{arXiv:2602.02276}. Qwen3.5-397B-A17B model card. Gemma Team, ``Gemma 4,'' \href{https://arxiv.org/abs/2607.02770}{arXiv:2607.02770}.}
\end{frame}

\begin{frame}{Vision Early: Two Independent Lines of Evidence}

    {\footnotesize \textbf{Terms.} \emph{Early fusion}: image patches enter the first layer, with no vision encoder. \emph{Encoder-based}: a pre-trained vision encoder feeds the language model (Shukor et al.\ call it late fusion).\par}
    \vspace{0.2em}
    \begin{columns}[T,onlytextwidth]
        \column{0.49\textwidth}
        \textbf{A vendor ablation (Kimi K2.5)}
        \begin{itemize}\small\itemsep0.2em
            \item With a fixed total budget of vision and text tokens, adding vision \textbf{early in training, at only 10\%} of the mix, beats adding it \textbf{late at 50\%}.
            \item The ratio matters little. The timing matters.
            \item A bonus: reinforcement learning on \emph{vision} tasks also raised \emph{text} scores, by about 2 points on GPQA-Diamond.
        \end{itemize}
        \column{0.49\textwidth}
        \textbf{An academic scaling study (Apple)}
        \begin{itemize}\small\itemsep0.2em
            \item 457 models, trained from scratch. At equal compute, early fusion is \textbf{on par} with the encoder-based design, and slightly better at small budgets.
            \item Both follow almost the same scaling exponent as text-only models:
        \end{itemize}
        \vspace{0.2em}
        \centering
        \footnotesize
        \begin{adjustbox}{max width=\linewidth}\begin{tabular}{lc}
            \toprule
            & $L\propto C^{\,x}$ \\
            \midrule
            Early fusion & $-0.0492$ \\
            Encoder-based & $-0.0494$ \\
            GPT-3 (text) & $-0.048$ \\
            \bottomrule
        \end{tabular}\end{adjustbox}
    \end{columns}
    \vspace{0.3em}
    \takeaway{An encoder's real advantage is \emph{reuse} of pre-trained parts. Enough native training shrinks it.}
    \src{Kimi Team, \href{https://arxiv.org/abs/2602.02276}{arXiv:2602.02276}, Tables 1 and 2. Shukor et al., ``Scaling Laws for Native Multimodal Models,'' ICCV 2025, \href{https://arxiv.org/abs/2504.07951}{arXiv:2504.07951}.}
\end{frame}

\begin{frame}{Sharing Capacity Between Modalities}
    \begin{itemize}\itemsep0.36em
        \item If one backbone handles every modality, should some weights be reserved for each?
        \item \textbf{Mixture-of-Transformers (Meta).} Give each modality its own feed-forward, attention and normalisation weights, but keep \emph{one global self-attention} over the whole sequence. Routing is by modality, with no learned router. It matches a dense Chameleon-style model at \textbf{55.8\%} of the FLOPs.
        \item \textbf{Learned MoE (Apple).} Let the router decide. Routing that ignores modality labels beats routing that uses them. Experts still specialise by modality on their own, mostly in the first and last layers.
        \item These agree more than they differ:
    \end{itemize}
    \vspace{0.07em}
    \takeaway{\textbf{Specialise the capacity. Keep the attention global.} It is the same division of labour as MoE in a text model: experts differ, mixing is shared.}
    \src{Liang et al., ``Mixture-of-Transformers,'' \href{https://arxiv.org/abs/2411.04996}{arXiv:2411.04996}. Shukor et al., \href{https://arxiv.org/abs/2504.07951}{arXiv:2504.07951}.}
\end{frame}

\begin{frame}{Fewer Visual Tokens, Lighter Encoders}
    \begin{itemize}\itemsep0.55em
        \item An image can cost over a thousand tokens. A video costs that per frame. With long context expensive, \textbf{visual-token compression} became a design axis of its own.
        \begin{itemize}\itemsep0.25em
            \item \textbf{Budgets:} Gemma 4 lets the caller choose 70 to 1,120 tokens per image.
            \item \textbf{Temporal pooling:} Kimi K2.5 averages groups of four video frames at the patch level, and fits about 4$\times$ longer video.
            \item \textbf{Optical compression:} DeepSeek-OCR reads a page of text from an image with about 97\% precision at under 10$\times$ fewer tokens than the text itself. At 20$\times$ it falls to about 60\%.
        \end{itemize}
        \item \textbf{Does the encoder have to exist?} Gemma 4's 12B model has none. Raw 48$\times$48 pixel patches go through a single linear projection into the LLM. Audio chunks are projected the same way.
    \end{itemize}
    \vspace{0.2em}
    \caveat{No published ablation yet compares encoder-free and encoder-based models at production scale. The other Gemma 4 sizes keep a ViT.}
    \src{Gemma Team, \href{https://arxiv.org/abs/2607.02770}{arXiv:2607.02770}. Kimi Team, \href{https://arxiv.org/abs/2602.02276}{arXiv:2602.02276}. Wei et al., ``DeepSeek-OCR: Contexts Optical Compression,'' \href{https://arxiv.org/abs/2510.18234}{arXiv:2510.18234}.}
\end{frame}

\begin{frame}{One Model That Understands and Generates}
    \begin{center}
    \small
    \renewcommand{\arraystretch}{1.0}
    \begin{adjustbox}{max width=\linewidth}\begin{tabular}{>{\raggedright\arraybackslash}p{4.3cm}>{\raggedright\arraybackslash}p{3.4cm}cc}
        \toprule
        \textbf{How images are generated} & \textbf{Example} & \textbf{GenEval} $\uparrow$ & \textbf{MMMU} $\uparrow$ \\
        \midrule
        Next-token prediction over image tokens & Emu3.5 (34B) & 0.86 & n/a \\
        Autoregressive text, flow-matching images & BAGEL (14B) & 0.82 & 55.3 \\
        Masked discrete diffusion for everything & LLaDA2.0-Uni (16B) & 0.89 & 50.1 \\
        Multimodal LLM with a diffusion head & InternVL-U (4B) & 0.85 & 54.7 \\
        \bottomrule
    \end{tabular}\end{adjustbox}
    \end{center}
    \begin{itemize}\itemsep0.13em
        \item \textbf{No design has won.} Four very different mechanisms land within a few points of each other. Sizes and settings differ, so read the numbers as indicative.
        \item \textbf{Do the two tasks help each other?} A 2026 analysis says yes at the level of representations, and no if both are forced through identical computation. Share the knowledge. Decouple the pathway.
    \end{itemize}
    \src{Scores as reported in: ByteDance Seed, ``BAGEL,'' \href{https://arxiv.org/abs/2505.14683}{arXiv:2505.14683}. \href{https://arxiv.org/abs/2603.09877}{arXiv:2603.09877} (InternVL-U). \href{https://arxiv.org/abs/2604.20796}{arXiv:2604.20796} (LLaDA2.0-Uni). \href{https://arxiv.org/abs/2601.02204}{arXiv:2601.02204} (NextFlow, for Emu3.5). \href{https://arxiv.org/abs/2609.01607}{arXiv:2609.01607}.}
\end{frame}

\begin{frame}{Multimodal: What to Remember}
    \begin{itemize}\itemsep0.6em
        \item The flagship of 2026 is \textbf{one natively multimodal model}, built on the same hybrid and MoE backbones as text models.
        \item \textbf{Bring vision in early.} Adding it early in training beats adding it late, and an encoder-free model matches an encoder-based one at equal compute.
        \item \textbf{Specialise capacity, share attention.} Modality-specific experts with global mixing.
        \item \textbf{Visual tokens are a budget} to be compressed and chosen, like context length.
        \item \textbf{Unified generation is unsettled.} Autoregression, flow matching and discrete diffusion all remain in play.
    \end{itemize}
    \vspace{0.3em}
    \caveat{Closed labs disclose almost no architectural detail. What we can teach rigorously comes from open-weight reports.}
\end{frame}
